\chapter{The trust boundary: assumptions, axiom footprint, inherited sorries, findings, coverage, what is open}
\label{chap:trust}

\section*{At a glance}
\begin{itemize}
\item \lead{Goal} State everything the eleven chapters before this one assume, inherit, or leave
open, and measure it against the tree.
\item \lead{Main results} The capstone carries fourteen binders, three bundles ---
\code{ErasureSpec} (7 fields), \code{EraserAsks} (4), \code{UpstreamAsks} (4) --- plus
\code{hbridge}, an implication to a two-field \code{ErasureBridge} no theorem inhabits.
Thirty-three nodes below cover them, a measured 33-name axiom footprint
(Assumption~\ref{asm:lean-reflection-axioms}), and seventeen inherited lean4lean \code{sorry}
sites (Assumption~\ref{asm:lean4lean-trust}).
\item \lead{Status} \stAssumed{} on every class-C/D/upstream node; \stInherited{} on the
lean4lean and reflection-axiom nodes; \stOpen{} on the two \code{ErasureBridge} fields and on
\code{UpstreamAsks.constOriginExcludes}, which is refuted. \hl{No environment satisfies any rung's
hypotheses} (Section~\ref{sec:trust:vacuity}).
\item \lead{Lean files} \code{ErasureSpec.lean}, \code{Upstream.lean}, \code{Origin.lean},
\code{Capstone.lean}, \code{test/Ledger.lean}, \code{test/Vacuity.lean}; \code{doc/trust.md},
\code{doc/coverage.md}, \code{doc/upstream-asks.md}, \code{doc/rework/03-DEV-FIX.md},
\code{doc/rework/07-STATUS.md}, \code{doc/rework/12-REPAIRS-W9.md}.
\end{itemize}

\noindent\lead{How to read the nodes below} A node carries a Lean field's name when one exists,
and never a proof: the dependency graph renders it stated-but-unproved, and every dependent
inherits that colour. A node with no Lean counterpart names the definition of the predicate it
states instead.

\par\smallskip\noindent
\small
\begin{tabular}{lllc}
\toprule
\textbf{assumption} & \textbf{class} & \textbf{stands at} & \textbf{status} \\
\midrule
\ref{asm:lean4lean-trust} & inherited & every source-term result & \stInherited{} \\
\ref{asm:lean-reflection-axioms} & inherited & capstone, 8 rungs, \code{oracle\_sound\_of\_run} & \stInherited{} \\
\ref{asm:EraserAsks-passes_monotone} & C & every rung & \stAssumed{} \\
\ref{asm:EraserAsks-passes_sound} & C & every rung & \stAssumed{} \\
\ref{asm:EraserAsks-oracle_false_refl} & C & every rung & \stAssumed{} \\
\ref{asm:EraserAsks-kernel_ind_head_true} & C & every rung & \stAssumed{} \\
\ref{asm:ErasureSpec-env_connect} & D & every rung & \stAssumed{} \\
\ref{asm:ErasureSpec-lookup_adequate} & D & every rung & \stAssumed{} \\
\ref{asm:ErasureSpec-fresh_names} & D & every rung & \stAssumed{} \\
\ref{asm:ErasureSpec-oracle_refl} & D & every rung & \stAssumed{} \\
\ref{asm:ErasureSpec-decl_adequate} & D & every rung & \stAssumed{} \\
\ref{asm:ErasureSpec-prim_monotone} & D & every rung & \stAssumed{} \\
\ref{asm:ErasureSpec-block_adequate} & D & every rung & \stAssumed{} \\
\ref{asm:UpstreamAsks-constsOrigin} & upstream ask 2 & every rung & \stAssumed{} \\
\ref{asm:UpstreamAsks-constOriginExcludes} & upstream ask 2 & every rung & \stOpen{} (refuted) \\
\ref{asm:UpstreamAsks-mkAppsInv} & upstream ask 9 & every rung & \stAssumed{} \\
\ref{asm:UpstreamAsks-indSpineInj} & upstream ask 10 & every rung & \stAssumed{} \\
\ref{asm:ErasureBridge-erasesEnv} & -- & capstone, every rung & \stOpen{} \\
\ref{asm:ErasureBridge-lowerEnv} & -- & capstone, every rung & \stOpen{} \\
\ref{asm:hargReach} & -- & G8 only (\code{hbody}) & \stAssumed{} \\
\ref{asm:htbl} & D & capstone, every rung & \stAssumed{} \\
\ref{asm:hsafe} & D & capstone, every rung & \stAssumed{} \\
\ref{asm:hblk} & D & capstone, every rung & \stAssumed{} \\
\ref{asm:hcb} & -- & capstone, every rung (\stChecked{} at G1) & \stAssumed{} \\
\ref{asm:hprep} & D & capstone, every rung & \stAssumed{} \\
\ref{asm:hrun} & D & capstone, every rung & \stAssumed{} \\
\ref{asm:hev} & C & capstone, every rung (\stChecked{} at G5) & \stAssumed{} \\
\ref{asm:hvwt} & C & capstone, every rung & \stAssumed{} \\
\ref{asm:hty} & C & capstone, every rung & \stAssumed{} \\
\ref{asm:hfo} & C & capstone, every rung & \stAssumed{} \\
\bottomrule
\end{tabular}
\normalsize

\section{The five trust classes}
\label{sec:trust:classes}

\code{doc/trust.md} grades every assumption on one scale, used by name in Chapters~\ref{chap:erases},
\ref{chap:eraser}, \ref{chap:refinement}, \ref{chap:coldstart} and \ref{chap:capstone}.

\par\smallskip\noindent
\begin{tabular}{lp{0.87\linewidth}}
\toprule
\textbf{class} & \textbf{meaning} \\
\midrule
A & proved here, footprint within \code{propext}, \code{Classical.choice}, \code{Quot.sound} \\
B & proved here, footprint additionally contains \code{sorryAx} inherited from lean4lean --- \eg\
\code{SEval.defeq} \\
C & a hypothesis of a stated theorem, visible in the statement \\
D & a permanent named binder, mechanised outside Lean \\
E & a scope restriction or a fact about a consumer, stated and not formalised \\
\bottomrule
\end{tabular}

\begin{itemize}
\item \lead{Pinned configuration} \code{ConfigPinned}: \code{csimp = false}, \code{extern =
.preferLogical}, \code{nat = .peano}, \code{remove\_irrel\_constr\_args = false},
\code{auto\_inline\_typeclass\_dispatch = false}.
\item \hl{Not the shipping default}: a bare \code{\#erase} sets \code{csimp} on, \code{extern :=
.preferAxiom} and machine \code{Nat} --- three of five fields differ, so a default invocation is
outside every theorem here.
\end{itemize}

The class-E rows are scope restrictions, stated and not formalised:

\par\smallskip\noindent
\small
\begin{tabular}{lp{0.76\linewidth}}
\toprule
\textbf{id} & \textbf{restriction} \\
\midrule
N1 & the \code{@[csimp]} replacement walk \\
N2 & \code{@[extern]} axiomatisation: fails \code{NoBodylessRefs} (Definition~\ref{def:NoBodylessRefs}), outside the domain \\
N3 & machine \code{Nat} \\
N4 & the constructor argument mask \\
N5 & auto-inlining and the \code{.ast.inlinings} sidecar \\
N7 & termination: no source strong-normalisation result; the capstone keeps [JACM]'s conditional form \\
N10 & serialisation \\
N11 & the remaining sidecars \\
N14 & size \\
N16 & \code{Quot} \\
N18 & elimination out of \code{Prop} \\
max-free levels & every tabled body lies in the \code{max}-free fragment (\code{TableSafe.noMaxLevels}), the fragment erasure transports along under level instantiation \\
-- & well-founded definitions: compiler bodies differ from kernel bodies; only propositional instances covered \\
-- & peregrine's own untyped-transform precondition is \code{Admitted}; \code{validate} checks no expandedness \\
-- & MetaRocq's shipped \code{firstorder\_ind} is \code{false} on \code{nat}; this development cites [JACM] rather than transcribing it \\
-- & no Rocq-side copy of $\Erases$ or $\Lower$ exists; \code{doc/rules-Erases.md}/\code{doc/rules-Lower.md} anchor instead \\
-- & F-PRODUCT: \code{auto\_inline\_typeclass\_dispatch}, off by default and pinned \code{false} \\
\bottomrule
\end{tabular}
\normalsize

\section{Inherited trust}
\label{sec:trust:inherited}

\begin{assumption}[The lean4lean trust boundary: seventeen inherited sorry sites]
\label{asm:lean4lean-trust}
\uses{def:lean4lean-grounding}
\stInherited{} \lead{In short} Every result here about \emph{source} terms is relative to
lean4lean's metatheory, and the pinned fork is not \code{sorry}-free.
\begin{itemize}
\item \lead{Pin} \code{8cc17a54c41c5e364596f8e4039e1f2094523a88} (the \code{trproj} branch of a
fork adding iota-reduction support, \code{TrProj} and the well-formed-environment origin
lemmas), recorded in \code{lakefile.toml}, both fields of \code{lake-manifest.json} and
\code{test/lean4lean-sorries.expected}; re-enumerated every CI run by
\code{scripts/lean4lean-sorries.sh}.
\item \lead{Seventeen sites}, by file, and whether they reach this tree.
\item This repository itself has \textbf{no} \code{sorry}: every \code{sorry} warning a build
reports is lean4lean's.
\end{itemize}
\par\smallskip\noindent
\small
\begin{tabular}{p{0.23\linewidth}lp{0.50\linewidth}l}
\toprule
\textbf{file} & \textbf{count} & \textbf{names} & \textbf{reaches?} \\
\midrule
ChurchRosser.lean & 2 & \code{VEnv.NormalEq.parRed} & yes \\
EnvLemmas.lean & 1 & \code{VEnv.WF.patsStrong} (fork-authored) & yes \\
Injectivity.lean & 4 & \code{sort\_inv}, \code{forallE\_inv\_stratified}, \code{sort\_forallE\_inv}, \code{const\_arity\_inv} & yes \\
UniqueTyping.lean & 1 & \code{IsDefEqU.weakN\_iff} & yes \\
Verify/Environment.lean & 1 & \code{addDecl.WF} (\code{inductDecl} case) & no \\
executable checker & 6 & \code{inferProj.WF\_struct}, \code{inferProj.WF}, \code{tryEtaStructCore.WF}, \code{isDefEqUnitLike.WF}, \code{reduceProjCore.WF}, \code{reduceRecursor.WF} & no \\
Verify/Typing/Lemmas.lean & 2 & \code{TrProj.weak'\_inv}, \code{TrProj.uniq} & no \\
\bottomrule
\end{tabular}
\normalsize
\end{assumption}

\begin{remark}
\lead{Caveat} Only a re-pin, inspected by fixture diff, discharges this. Seven named declarations
reach this tree: six through \code{TrExprS.uniq} and \code{IsDefEq.uniqU}, and
\code{IsDefEqU.const\_arity\_inv} by direct citation from \code{not\_erasable\_of\_informative},
\code{indSpine\_ne\_forallE} and \code{FirstOrderInd.notSortNotPi}. \code{step\_proj} carries
exactly the iota arm's \code{sorryAx} roots: it \emph{consumes} a \code{TrProj} witness rather
than building one, an inhabitation gap, not an axiom one.
\end{remark}

\begin{assumption}[Lean's reflection axioms: the price of the oracle reroute]
\label{asm:lean-reflection-axioms}
\uses{asm:ErasureSpec-oracle_refl}
\stInherited{} \lead{In short} The capstone and all eight rungs carry a 33-name axiom footprint:
four standard names plus 29 kernel-reflection axioms, the price of routing the oracle through
lean4lean's executable checker instead of assuming its soundness.
\begin{itemize}
\item \lead{Class} inherited: Lean's and lean4lean's own modelling of opaque, \code{@[extern]}
primitives, not proof holes.
\item \lead{Stands at} the capstone, all eight rungs (via \code{erasure\_bridge\_of\_run}'s oracle
step) and \code{ErasureSpec.oracle\_sound\_of\_run}. Four of the 29 also reach the bridge
\emph{outside} the oracle, as the footprint of \code{BridgeInv.mkLocalDecl}/\code{mkLetDecl}; a
six-name subset is what the delta arm adds to \code{erases\_correct} (Section~\ref{sec:trust:ledger}).
\item \lead{Would be discharged by} reverting the kernel reroute (about forty lines, five hunks)
and keeping the oracle's reflection clause as a class-D binder --- at the price of assuming what is
proved.
\item \lead{Caveat} \code{Level.hasParam\_eq} is the bridge from Lean's cached level bitfield to a
structural mirror, so every kernel-grade statement about level parameters inherits it.
\end{itemize}
\par\smallskip\noindent
\small
\begin{tabular}{p{0.20\linewidth}lp{0.45\linewidth}p{0.17\linewidth}}
\toprule
\textbf{group} & \textbf{count} & \textbf{names} & \textbf{site} \\
\midrule
pointer equality & 2 & \code{ptrEqExpr\_eq}, \code{ptrEqConstantInfo\_eq} & PtrEq.lean \\
\code{Lean.Expr} reflection & 14 & \code{abstract\_eq}, \code{instantiate\_eq}, \code{replace\_eq}, ... & Verify/Axioms.lean \\
\code{Lean.Level} & 5 & \code{hasMVar\_eq}, \code{normalize\_eq}, ... & Verify/Axioms.lean \\
persistent containers & 4 & \code{PersistentArray.toList'\_push}, \code{PersistentHashMap.WF.find?\_eq}, ... & Verify/Axioms.lean \\
\code{Std.TreeMap} & 1 & \code{all\_eq\_all\_toList} & Verify/Axioms.lean \\
syntax & 1 & \code{Lean.Syntax.structEq\_eq} & Verify/Axioms.lean \\
\code{bv\_decide} certificates & 2 & \code{Expr.mkData\_flags}, \code{Expr.Data.looseBVarRange\_le} & Verify/Expr.lean \\
\bottomrule
\end{tabular}
\normalsize
\end{assumption}

\section{Class C: what the eraser asks of itself}
\label{sec:trust:eraser}

\code{EraserAsks lenv env gw} (Definition~\ref{def:EraserAsks}) bundles \textbf{four} obligations
about code \emph{this repository owns} --- the preparation passes and the relevance oracle. It
binds \code{erasure\_bridge\_of\_run}, the capstone and every rung; \textbf{no rung discharges any
field}. The bundle takes no level scope: each field speaks of the scope its own run is made at.
Two fields yield one derived theorem, \code{EraserAsks.oracle\_informative}.

\begin{assumption}[The preparation passes only advance the name generator]
\label{asm:EraserAsks-passes_monotone}
\lean{LeanToLambdaBox.EraserAsks.passes_monotone}
\leanok
\uses{def:preparePasses, def:Erasure-EraseM}
\stAssumed{} \lead{In short} Every \code{CoreM} pass of \code{preparePasses}, lifted into the
erasure monad, leaves the ambient name generator \hl{at least as advanced}.
\begin{itemize}
\item \lead{Given} the three passes --- \code{replaceUnsafeRecNames}, \code{macroInline},
\code{inlineMatchers} --- quantified over membership, not the call sequence: preparation calls
them four times, \code{macroInline} twice; the \code{@[csimp]} walk is out of scope.
\item \lead{Class} C; \lead{Stands at} every rung, discharged nowhere.
\item \lead{Would be discharged by} unfolding \code{Lean.Core.transform}'s generator discipline.
\end{itemize}
\end{assumption}

\begin{assumption}[The preparation passes preserve source evaluation, at any spine]
\label{asm:EraserAsks-passes_sound}
\lean{LeanToLambdaBox.EraserAsks.passes_sound}
\leanok
\uses{def:preparePasses, def:SEval}
\stAssumed{} \lead{In short} Each pass's run from $e$ to $e'$ preserves $\SEval$ at \hl{every
argument spine}: $\SEval\ (\mkApps\ e\ args)\ v \Rightarrow \SEval\ (\mkApps\ e'\ args)\ v$.
\begin{itemize}
\item \lead{Given} the spine is universally quantified because the capstone reads its observable
at $\mkApps\ e\ args$ while a pass is a whole-tree walk, and $f(\mkApps\ e\ args) \neq \mkApps\
(f\,e)\ args$.
\item \lead{Class} C; \lead{Stands at} every rung, discharged nowhere alone --- but its
composition over the four calls, \code{prepare\_sound} (Theorem~\ref{thm:prepare_sound}), \emph{is}
proved, and is what the capstone spends.
\item \lead{Would be discharged by} a delta-expansion argument for \code{macroInline} and a
renaming argument for \code{replaceUnsafeRecNames}.
\end{itemize}
\end{assumption}

\begin{assumption}[A false oracle verdict reflects the pure kernel run]
\label{asm:EraserAsks-oracle_false_refl}
\lean{LeanToLambdaBox.EraserAsks.oracle_false_refl}
\leanok
\uses{def:Erasure-isErasable, def:isErasable}
\stAssumed{} \lead{In short} If the shipping oracle \code{Erasure.isErasable}, lifted into the
erasure monad, returns \code{false}, the \hl{pure verified checker} did not answer \code{true} at
the same local context and level scope.
\begin{itemize}
\item \lead{Given} both runs are read at \code{ctx.lparams}, the one scope
\code{Erasure.isErasable} is called at; no reader's scope enters the statement. The pure run is
\code{LeanToLambdaBox.isErasable} via \code{Lean4Lean.TypeChecker.M.run} at \code{lenv}'s kernel
environment.
\item \lead{Class} C; \lead{Stands at} every rung, discharged nowhere.
\item Near-definitional; a \code{MetaM} reflection lemma for the success arm would discharge it.
\end{itemize}
\end{assumption}

\begin{assumption}[At an inductive-type head the kernel oracle answers true]
\label{asm:EraserAsks-kernel_ind_head_true}
\lean{LeanToLambdaBox.EraserAsks.kernel_ind_head_true}
\leanok
\uses{def:IndInfo, def:isErasable, def:lean4lean-grounding}
\stAssumed{} \lead{In short} At an inductive-type head with a well-formed translated context, the
pure checker answers \code{true} --- the oracle's \hl{completeness} at the one shape the fragment
cannot exclude.
\begin{itemize}
\item \lead{Given} a level scope \code{lps} quantified by the field --- the run's own, the one
\code{visitMutual} installs from the declaration being entered, not a reader's;
\item $e$'s head is a constant the model declares an inductive type former, and $e$ translates
($\TrExprS$) at \code{lps} and a well-formed modelled local context whose free variables
\code{kernelNGen}, this repository's name for the checker's initial generator, reserves.
\item \lead{Class} C; \lead{Stands at} every rung, discharged nowhere, and \textbf{not a
theorem}: \code{isArityCheck} walks the head's type under a fixed budget, so a \emph{reduced}
telescope longer than that budget --- or any other kernel error inside the check --- leaves the
walk short of the closing sort and routes the verdict to the \code{Meta} fallback, of which only
soundness is assumed (the second disjunct of Assumption~\ref{asm:ErasureSpec-oracle_refl}).
\item \lead{Measured} \code{.ok true} at 2{,}940 of 2{,}940 inductive type formers of the
elaboration environment.
\item \lead{Would be discharged by} three executable-shape lemmas lean4lean lacks, together with
a bound on the reduced telescope that a constant budget does not supply.
\end{itemize}
\end{assumption}

\begin{remark}
\lead{Why} \code{EraserAsks} carries no key-distinctness field. A mutual block's $\lbox$ keys are
distinct because \code{visitMutual} refuses a block whose mapped kernames repeat, and
\code{Erasure.run\_rec\_exit\_nodup} reads that guard off a successful run; the distinctness the
install site consumes is a run conclusion, not a standing claim about
\code{Lean.Compiler.LCNF.getDeclInfo?}, which is false (F-UNSAFEREC).
\end{remark}

\section{Class D: the specification of Lean's primitives}
\label{sec:trust:spec}

\code{ErasureSpec lenv env Us gw} (Definition~\ref{def:ErasureSpec}) specifies the ambient
primitives the erasure runs on, in \textbf{seven} fields. \lead{Class D throughout}: each field is
a permanent binder about an object no Lean term denotes (\code{IO.RealWorld} opaque,
\code{Lean.Environment} has a private constructor); three fields are partly \emph{earned} (below).
The capstone and every rung bind it at every level scope, \code{P :\ forall Us, ErasureSpec lenv env Us gw}.

\begin{assumption}[The elaboration environment is modelled by the specification environment]
\label{asm:ErasureSpec-env_connect}
\lean{LeanToLambdaBox.ErasureSpec.env_connect}
\leanok
\uses{def:lean4lean-grounding}
\stAssumed{} \lead{In short} \code{env} is the \code{.safe} member of a family of model
environments, well formed for the kernel environment underlying \code{lenv}.
\begin{itemize}
\item \lead{Given} the kernel-environment conversion is explicit, matching lean4lean's own
well-formedness (at \code{Lean.Kernel.Environment}) and the shipping oracle's conversion.
\item \lead{Class} D; \lead{Stands at} every rung, discharged nowhere; \lead{Would be discharged
by} upstream ask 4.
\item \lead{Buys} \code{ErasureSpec.envWF} (Theorem~\ref{thm:ErasureSpec-envWF}): earned, not
assumed, model well-formedness; the capstone spends it twice.
\end{itemize}
\end{assumption}

\begin{assumption}[Environment lookups report what the elaboration environment holds]
\label{asm:ErasureSpec-lookup_adequate}
\lean{LeanToLambdaBox.ErasureSpec.lookup_adequate}
\leanok
\uses{def:LookupAdequate, def:CasesOnShape}
\stAssumed{} \lead{In short} The four environment queries a run makes report what \code{lenv}
holds and \hl{none advances the name generator} --- one application of \code{LookupAdequate}.
\begin{itemize}
\item \code{constInfo}: \code{getConstInfo} returns \code{lenv}'s own declaration.
\item \code{declInfo}: the name lies in \code{getDeclInfo?}'s block (guard: not an
\code{\_unsafe\_rec} companion); \code{none} only for an unknown name.
\item \code{ctorArity}: \code{getCtorArity?} answers exactly for \code{lenv}'s constructors ---
the refinement's fourth motive needs the negative direction.
\item \code{casesInfo}: \code{getCasesInfo?} answers exactly for the \code{casesOn} constants, at
metadata agreeing with the declared block through \code{CasesInfoAgreesK}: seven clauses, among
them \code{indName} --- the eliminated inductive is the one the alternatives are read against, not
the head's name prefix --- and \code{altCtor}, slot $j$ is constructor $j$'s and never the
catch-all shape. Measured at 3{,}161 of 3{,}161 \code{casesOn} constants, 0 failures.
\item \lead{Class} D; \lead{Stands at} every rung, discharged nowhere.
\end{itemize}
\end{assumption}

\begin{assumption}[Fresh free variables are fresh]
\label{asm:ErasureSpec-fresh_names}
\lean{LeanToLambdaBox.ErasureSpec.fresh_names}
\leanok
\stAssumed{} \lead{In short} \code{Lean.mkFreshFVarId}, in the erasure monad, returns an
identifier \hl{not yet handed out}, reserves it, only advances the generator, and returns one the
kernel's own generator reserves too.
\begin{itemize}
\item \lead{Given} a four-way conjunction; the kernel generator is \code{kernelNGen}, an
\code{abbrev} this repository declares in its own namespace for the checker's initial generator.
\item \lead{Class} D; \lead{Stands at} every rung, discharged nowhere --- the generator is read
from the opaque \code{IO.RealWorld}, so no Lean proof exists in the present formulation.
\end{itemize}
\end{assumption}

\begin{assumption}[The monadic oracle run reflects the kernel predicate]
\label{asm:ErasureSpec-oracle_refl}
\lean{LeanToLambdaBox.ErasureSpec.oracle_refl}
\leanok
\uses{def:Erasure-isErasable, def:isErasable, def:Oracle-MetaSound}
\stAssumed{} \lead{In short} A \code{true} verdict of \code{Erasure.isErasable} taken at scope
\code{Us} either reflects a successful pure checker run, or came from the
\code{isErasableMeta} fallback, \hl{assumed sound} (\code{Oracle.MetaSound}).
\begin{itemize}
\item \lead{Given} the run only advances the generator; the disjunction is read under
\code{ctx.lparams = Us}. The capstone takes the bundle at \emph{every} \code{Us}, so a run
below \code{visitMutual}'s re-entry spends it at the member's own level column. Irreducibly, the
impure \code{MetaM} plumbing reflects the pure kernel run it calls.
\item \lead{Class} D; \lead{Stands at} every rung, discharged nowhere as stated --- but the kernel
disjunct is \emph{spent, not assumed}: \code{ErasureSpec.oracle\_sound\_of\_run}
(Theorem~\ref{thm:ErasureSpec-oracle_sound_of_run}) composes it with
\code{Oracle.kernel\_isErasable\_sound} to conclude $\Erasable$ outright, bringing the 29
reflection axioms of Assumption~\ref{asm:lean-reflection-axioms} into the capstone's footprint.
\item \lead{Measured} 0 fallback hits in 139{,}196 constants.
\end{itemize}
\end{assumption}

\begin{assumption}[A declaration visible at safe is visible in the model]
\label{asm:ErasureSpec-decl_adequate}
\lean{LeanToLambdaBox.ErasureSpec.decl_adequate}
\leanok
\uses{def:lean4lean-grounding}
\stAssumed{} \lead{In short} Whenever \code{lenv.find?} answers with a \code{ConstantInfo} at
safety \code{.safe}, the model declares that name with a translated type, \hl{related by
\code{TrConstant}}.
\begin{itemize}
\item \lead{Given} in lean4lean's order \code{.safe} is the \emph{top} of
\code{DefinitionSafety}, so the statement's inequality (reading backwards) forces exactly
\code{.safe}.
\item \lead{Class} D; \lead{Stands at} every rung. \code{decl\_adequate\_of\_kernelFind}
\emph{proves} it from Assumption~\ref{asm:ErasureSpec-env_connect} for the \emph{kernel}
environment's own lookup; it cannot cross \code{Lean.Environment.find?} versus
\code{Lean.Kernel.Environment.find?}, two opaque primitives whose agreement is not a theorem here.
\item \lead{Would be discharged by} nothing upstream alone: ask 4 covers the kernel lookup, and
the two lookups' agreement stays a specification decision.
\end{itemize}
\end{assumption}

\begin{remark}
\lead{Caveat} This feeds Assumption~\ref{asm:hsafe}: its three safety clauses exist because
\code{decl\_adequate} is stated at safe declarations only, while the table's adequacy predicate
does not record the safety flag.
\end{remark}

\begin{definition}[Generator monotonicity of a MetaM computation]
\label{def:MetaGenMono}
\lean{LeanToLambdaBox.MetaGenMono}
\leanok
\lead{In short} A \code{Lean.MetaM} computation only advances the name generator, stated at the
shape such a computation is \hl{applied} in.
\begin{itemize}
\item \lead{Where} the applied shape takes a \code{Meta.Context}, a \code{Meta.State} reference, a
\code{Core.Context}, a core state reference and a world token;
\item at that shape \code{>>=} is the underlying \code{EST.bind} definitionally, so the property
composes through a \code{do} block (\code{MetaGenMono.bind}, \code{.pure}, the two \code{forIn}
combinators);
\item at the \code{Erasure.liftMetaM} shape it does not compose, because \code{MetaM.run'}
allocates a fresh \code{Meta.State} reference per call --- crossing between the two shapes is what
\code{PrimMonotone.liftMetaM} is for.
\end{itemize}
\end{definition}

\begin{definition}[The derived MetaM bound: no quantifier over arbitrary computations]
\label{def:PrimGenMono}
\lean{LeanToLambdaBox.PrimGenMono}
\leanok
\uses{def:PrimMonotone, def:MetaGenMono}
\lead{In short} The \code{MetaM} computations whose generator bound \hl{follows from}
\code{PrimMonotone} rather than being assumed of them.
\begin{itemize}
\item \lead{Where} $x$ is \code{PrimGenMono} when every \code{PrimMonotone gw} gives
\code{MetaGenMono gw x};
\item every occurrence is discharged by the \code{pure}, \code{>>=} and \code{forIn} combinators
on the three named primitives --- \code{isProof} and the two bounded telescopes;
\item it is what covers the two anonymous continuations the erasure passes to the telescopes, so
no clause of the bundle quantifies over an arbitrary \code{Lean.MetaM} computation.
\end{itemize}
\end{definition}

\begin{assumption}[The effect-only primitive calls only advance the generator]
\label{asm:ErasureSpec-prim_monotone}
\lean{LeanToLambdaBox.ErasureSpec.prim_monotone}
\leanok
\uses{def:PrimMonotone}
\stAssumed{} \lead{In short} The calls the erasure makes for effect alone only advance the name
generator --- one application of \code{PrimMonotone}, \hl{eight clauses, one per primitive}.
\begin{itemize}
\item \lead{Given} four named calls: \code{Lean.getEnv}, \code{Lean.logInfo},
\code{Lean.Meta.isInstance} and \code{Lean.Meta.inferType}, whose Pi-telescope additionally
matches the subject's lambda-telescope (\code{ForallMatchesLam}, vacuous unless the term side is a
lambda);
\item \code{Lean.Meta.isProof}, the proof test run on a constructor's fields and on a catch-all's
hypotheses, as a \code{MetaGenMono} clause (Definition~\ref{def:MetaGenMono});
\item \code{Lean.Meta.forallBoundedTelescope} and \code{Lean.Meta.lambdaBoundedTelescope}, each
\emph{compositionally}: the telescope opens binders of its own, so it advances the generator, and
no further than its continuation does;
\item \code{Erasure.liftMetaM}, which carries a bound from the applied \code{MetaM} shape to the
shape the run lemmas read;
\item \lead{Class} D; \lead{Stands at} every rung, discharged nowhere.
\end{itemize}
\end{assumption}

\begin{assumption}[The kernel's inductive blocks are the model's]
\label{asm:ErasureSpec-block_adequate}
\lean{LeanToLambdaBox.ErasureSpec.block_adequate}
\leanok
\uses{def:BlockAdequate, def:IndInfo, def:CasesOnShape}
\stAssumed{} \lead{In short} The kernel's inductive blocks, constructors and eliminators are the
model's, at the identifier \hl{\code{register\_inductive} mints} --- one application of the
eight-field \code{BlockAdequate}.
\begin{itemize}
\item \code{fwd}/\code{bwd}: \code{IndInfo} relates to a declared block with its kernel fields.
\item \code{ctor}/\code{ctorBwd}: \code{CtorOf} relates to a declared \code{ConstructorVal}.
\item \code{casesOnDecl}: the eliminator's \code{CasesOnShape} pins the block's own segmentation
--- parameters, discriminee, indices --- to the elaborator's.
\item \code{selfMem}, \code{selfName}: a declared inductive is a member of its own block, under
its own name.
\item \code{fields}: a declared inductive's constructor records are the kernel's own, at their
positions.
\item \lead{Caveat} the emitted propositional flag is \emph{not} a ninth clause: the half the
consumers spend is derived (Theorem~\ref{thm:ErasureSpec-propositionalInd_of_arity}).
\item \lead{Class} D; \lead{Stands at} every rung, discharged nowhere.
\end{itemize}
\end{assumption}

\begin{theorem}[A true arity verdict is propositionality in the model]
\label{thm:ErasureSpec-propositionalInd_of_arity}
\lean{LeanToLambdaBox.ErasureSpec.propositionalInd_of_arity}
\leanok
\uses{asm:ErasureSpec-decl_adequate}
\lead{In short} A \code{true} verdict of the eraser's own arity walk on a declared inductive's
type makes that inductive \hl{propositional in the model}.
\begin{itemize}
\item \lead{Given} \code{ErasureSpec lenv env Us gw};
\item \code{lenv} declares $I$ as an \code{inductInfo} at safety \code{.safe};
\item \code{Erasure.isPropositionalArity} answers \code{true} on its declared type;
\item \lead{Then} \code{PropositionalInd env I}: the declared arity's result level is zero at
every valuation.
\end{itemize}
\lead{Status} \stProved{}; conditional on \stAssumed{} \code{ErasureSpec}.
\end{theorem}
\begin{proof}
\leanok
\uses{asm:ErasureSpec-decl_adequate}
\lead{Idea} \code{decl\_adequate} translates the declared type, and the model-side result-sort
walk commutes with the syntactic one arm for arm, \code{let} and \code{mdata} included.
\lead{Steps} The converse is false: \code{arity\_of\_propositionalInd\_false} refutes it at an
arity whose final sort is a \code{let}'s own bound variable, so only the direction the iota and
projection arms spend is claimed.
\end{proof}

\section{The four upstream asks every rung depends on}
\label{sec:trust:upstream}

\code{UpstreamAsks env} (Definition~\ref{def:UpstreamAsks}) packages items 2, 9 and 10 of
\code{doc/upstream-asks.md} as four fields, one of them refuted. It binds \code{erases\_correct}
and eight further theorems, \code{erasure\_bridge\_of\_run}'s steps 3, 4 and 17, and the
accumulator bundle \code{AccAsks}, so \textbf{every rung carries it}. Item 6 is not a field: its
three consumers cite lean4lean's \code{IsDefEqU.const\_arity\_inv}, an inherited \code{sorry}
(Assumption~\ref{asm:lean4lean-trust}). Not measurable in the ledger: \code{\#print axioms}
reports a proved term's footprint, never a hypothesis.

\begin{assumption}[Ask 2: the uniqueness corollaries of a constant's origin]
\label{asm:UpstreamAsks-constsOrigin}
\lean{LeanToLambdaBox.UpstreamAsks.constsOrigin}
\leanok
\uses{def:CtorOf, def:IndInfo, def:CasesOnShape, def:IndDeclOf, def:IndBlockBelow}
\stAssumed{} \lead{In short} Five conjuncts fixing that a declared constant's origin data is
\hl{unique}.
\begin{itemize}
\item \lead{Given} a constructor belongs to one type at one index; \code{IndInfo} and
\code{CasesOnShape} are functional in their data; an \code{IndInfo} yields an
\code{IndDeclOf}; two blocks declared below \code{env} (\code{IndBlockBelow}) that declare the
same type former are the same block.
\item \lead{Class} upstream ask 2; \lead{Stands at} every rung, discharged nowhere here.
\item \lead{Not in the field} the classification conjunct, which \code{consts\_classified}
derives off lean4lean's \code{VEnv.WF'.consts\_origin}; and \code{CtorOf} excluding
\code{IndInfo}, which \code{CtorOf.not\_indInfo} proves.
\item \lead{Would be discharged by} reading the origin predicates off \code{env}'s own
declaration list (Section~\ref{sec:trust:open}), a downstream restatement, not a fork change.
\end{itemize}
\end{assumption}

\begin{assumption}[Ask 2: a plain constant is neither a constructor nor a type former]
\label{asm:UpstreamAsks-constOriginExcludes}
\lean{LeanToLambdaBox.UpstreamAsks.constOriginExcludes}
\leanok
\uses{def:ConstOrigin, def:CtorOf, def:IndInfo, lem:CtorOf-constOrigin}
\stOpen{} \lead{In short} Every \code{ConstOrigin} constant is neither a constructor nor an
inductive type name --- \hl{refuted at every environment declaring an inductive block}.
\begin{itemize}
\item \lead{Refuted} Lemma~\ref{lem:CtorOf-constOrigin} makes \code{Nat.zero} a
\code{ConstOrigin} at the fixture \code{natEnv}; \code{test/Vacuity.lean}'s
\code{not\_upstreamAsks\_natEnv} derives \code{False} from this field there.
\item \lead{Consequence} \code{UpstreamAsks} is inhabited at no environment that declares an
inductive block (Section~\ref{sec:trust:vacuity}).
\item \lead{Class} upstream ask 2, but no fork change answers it: the defect is
\code{ConstOrigin}'s reading, where MetaRocq's \code{declared\_constant} reads
\code{lookup\_env}. Consumers: \code{constOrigin\_not\_ctorOf}, \code{constOrigin\_not\_indInfo}.
\item \lead{Would be discharged by} the declaration-list restatement of
Section~\ref{sec:trust:open}, after which exclusion follows from lean4lean's
\code{consts\_origin}.
\end{itemize}
\end{assumption}

\begin{assumption}[Ask 9: spine typing inversion]
\label{asm:UpstreamAsks-mkAppsInv}
\lean{LeanToLambdaBox.UpstreamAsks.mkAppsInv}
\leanok
\uses{def:IndBlockBelow, def:lean4lean-grounding, asm:lean4lean-trust}
\stAssumed{} \lead{In short} A typed application spine's head has a type, and the
argument-by-argument \code{Peel} relation holds down the spine, under the model's \hl{strong-order
well-formedness}.
\begin{itemize}
\item \lead{Given} the strong order is an \emph{explicit} hypothesis, so the ask is
\code{sorryAx}-free like the single-application inversion. Consumers: \code{indSpine\_not\_prop},
\code{elim\_major}, \code{ctor\_saturated}, \code{fOFields\_of\_asks}.
\item \lead{Class} upstream ask 9; \lead{Stands at} every rung, discharged nowhere here;
\lead{Would be discharged by} a pin bump landing the fork's \code{HasType.mkApps\_inv}.
\item \lead{Caveat} it does \emph{not} buy sorry-freedom: the strong order's only introduction
rests on \code{VEnv.WF.patsStrong}, \code{sorry} at the pin.
\end{itemize}
\end{assumption}

\begin{assumption}[Ask 10: definitionally equal inductive spines have the same head]
\label{asm:UpstreamAsks-indSpineInj}
\lean{LeanToLambdaBox.UpstreamAsks.indSpineInj}
\leanok
\uses{def:IndDeclOf, def:lean4lean-grounding, asm:lean4lean-trust}
\stAssumed{} \lead{In short} Two definitionally equal application spines headed by
\emph{inductively declared} type formers have \hl{the same head}.
\begin{itemize}
\item \lead{Given} the \code{IndDeclOf} premises are essential: without them the statement is
false, refuted by a plain definition whose body is an inductive type.
\item \lead{Class} upstream ask 10; \lead{Stands at} every rung, discharged nowhere here.
Consumers: \code{ctor\_saturated}, \code{fOFields\_of\_asks}.
\item \lead{Would be discharged by} a pin bump landing the fork's \code{IsDefEqU.indSpine\_inj},
beside the three injectivity \code{sorry}s and Church--Rosser.
\end{itemize}
\end{assumption}

\section{No environment satisfies a rung's hypotheses}
\label{sec:trust:vacuity}

The refutation of Assumption~\ref{asm:UpstreamAsks-constOriginExcludes} reaches every rung and
the capstone's observable. \code{doc/rework/07-STATUS.md} sections 1 and 4 carry the detail.

\begin{itemize}
\item \lead{Every rung} binds \code{A :\ UpstreamAsks env} beside \code{hfo :\ FirstOrderInd env
Nat} and beside \code{P}, \code{htbl}, whose tabled \code{Nat} block gives \code{IndInfo env Nat}
(\code{indInfo\_of\_tabled}). \hl{Each pairing with \code{A} is \code{False}}, mechanised at the
three standard axioms.
\item \lead{The capstone} guards its observable by \code{FirstOrderInd env I} under the
conclusion's $\forall I$; under \code{A} it applies only at an environment declaring no inductive.
\item \lead{A false conclusion behind it} At \code{natEnv}, \code{Nat.zero} erases both by
\code{Erases.ctor} and by \code{Erases.const} to distinct $\lbox$ terms
(\code{test/Vacuity.lean}, \code{erases\_natZ\_two\_ways}), refuting the rungs' uniqueness conjunct;
\code{A} keeps it unreachable.
\item \lead{What the green rungs show} \code{green-check --all} is 8/8 and every rung elaborates:
evidence that a checked term meets the \emph{stated} interface, not that anything inhabits it.
The same holds for every accumulator step premised on \code{AccAsks}, which carries \code{A}.
\item \lead{Invisible to the ledger}: \code{\#print axioms} measures a proof, never whether its
hypotheses are satisfiable.
\end{itemize}

\section{The residual: the bridge's two environment fields}
\label{sec:trust:bridge}

Beyond the erasure half, the capstone assumes one named binder, \code{hbridge}: read at the state
the \code{visitExpr} run ends in, there is a specification environment $\Gamma_{\mathrm{spec}}$
adequate for that state, and \code{ErasureBridge} (Definition~\ref{def:ErasureBridge}, exactly
\textbf{two} fields) holds for every $\lbox$ term the prepared subject erases and lowers to.

\begin{itemize}
\item \lead{Proved interfaces} \code{bridgeEnv\_of\_regInv} (Theorem~\ref{thm:bridgeEnv_of_regInv})
derives both fields from the registration invariant \code{RegInvShape'} and saturation at the
run's final state. \code{bridgeEnv\_of\_regContent} (Theorem~\ref{thm:bridgeEnv_of_regContent}) composes the whole payload from \code{RegAcc},
\code{RegKeyed} and \code{ErasuresDeclared} at that state, plus the table's origin and
level-column clauses. Both measure \code{[propext, Classical.choice, Quot.sound]}.
\item \lead{Producers} \code{regKeyed\_of\_run} (Theorem~\ref{thm:regKeyed_of_run}) gives
\code{RegKeyed} at a run; \code{regInv\_registerInd\_run} and \code{regInv\_recConst\_step}
(Theorems~\ref{thm:regInv_registerInd_run} and~\ref{thm:regInv_recConst_step}) preserve
\code{RegAcc} across \emph{one} block registration and one recursive-constant registration.
\item \lead{The gap} \hl{nothing chains them across a whole run}: the accumulator step
\code{StepAcc6} is open, so \code{bridgeEnv\_of\_regContent} has no consumer and is not on the
capstone's proof path. \code{ErasuresDeclared} is refuted at a block member's sub-run.
\end{itemize}

\begin{assumption}[The specification environment erases the source environment]
\label{asm:ErasureBridge-erasesEnv}
\lean{LeanToLambdaBox.ErasureBridge.erasesEnv}
\leanok
\uses{def:ErasesEnv}
\stOpen{} \lead{In short} $\ErasesEnv$ holds of the specification environment along the erasure's
closure: every constant, block and axiom the erased term reaches is declared in
$\Gamma_{\mathrm{spec}}$ as the relation demands.
\begin{itemize}
\item \lead{Given} the relation reads a tabled body at \hl{the declaration's own level column},
which the bundle carries as a third environment parameter; that is where MetaRocq's
\code{erases\_constant\_body} reads it, and the instantiated form is derived at the delta step
rather than demanded here.
\item \lead{Stands at} the capstone and every rung, discharged nowhere:
\code{RegInvShape'.erasesEnv} derives it from the registration invariant at the run's final
state, but no theorem establishes that invariant for a whole run of the shipping eraser.
\item \lead{Would be discharged by} the accumulator induction (\code{StepAcc6}, then
\code{erasure\_bridge\_env}; \code{doc/rework/12-REPAIRS-W9.md}), after the declaration-list
restatement (Section~\ref{sec:trust:open}).
\end{itemize}
\end{assumption}

\begin{remark}
\lead{Caveat} This is the single most important caveat of the whole blueprint: the crown
theorem's environment half is assumed. No statement here may present the capstone, or any rung, as
unconditionally established.
\end{remark}

\begin{assumption}[The emitted environment is the lowered, pruned specification environment]
\label{asm:ErasureBridge-lowerEnv}
\lean{LeanToLambdaBox.ErasureBridge.lowerEnv}
\leanok
\uses{def:LowerEnv}
\stOpen{} \lead{In short} $\LowerEnv$ holds between the specification environment and the emitted
one: the emitted environment is the \hl{lowering image} of the specification environment, pruned
to what the program reaches.
\begin{itemize}
\item \lead{Given} its \code{defs} clause admits two shapes for a declared body: a $\Lower$ image,
or the eta-expansion of one lowered block's node, which is what the registration writes for a
recursive member.
\item \lead{Stands at} the capstone and every rung, discharged nowhere, for the reason
Assumption~\ref{asm:ErasureBridge-erasesEnv} is: derived from the same invariant at a saturated
registry, and nothing inhabits that invariant at a whole run.
\end{itemize}
\end{assumption}

\begin{remark}
\lead{Caveat} \code{bridgeEnv\_of\_regInv} and \code{bridgeEnv\_of\_regContent} are
\emph{proved}, at the three standard axioms and no \code{sorryAx}. The gap is the inhabitation of
their premises, not the implication.
\end{remark}

\begin{assumption}[G8 only: every erasure of the tabled body reaches Nat's block]
\label{asm:hargReach}
\uses{def:ReachableFrom, def:Erases}
\stAssumed{} \lead{In short} At G8 alone, as the binder \code{hbody}: for every specification
environment, every erasure of \code{benchArith}'s tabled body reaches the $\lbox$ key of
\hl{\code{Nat}'s inductive block}.
\begin{itemize}
\item \lead{Stands at} G8 only; stated inline in \code{green\_G8}'s signature.
\code{Green.g8\_hargReach} composes it into the per-argument $\ErasesEnv$ conjunct. G1--G7 apply
the observable clause at the empty spine, where that premise is vacuous.
\item \lead{Caveat} quantified over every $\Gamma_{\mathrm{spec}}$, because none is in scope at
the binder: strictly stronger than a reading at the run's own. Satisfiable and not refutable:
the body's literals erase through \code{Erases.lit} to peano towers naming \code{Nat}'s block.
\item \lead{Would be discharged by} a source-side argument about $\Erases$ at the reified body,
not by the registration invariant.
\end{itemize}
\end{assumption}

\section{The standalone capstone and rung binders}
\label{sec:trust:binders}

\begin{itemize}
\item \lead{Capstone's 14 explicit binders}, in order: \code{P}, \code{E}, \code{A}, \code{htbl},
\code{hsafe}, \code{hblk}, \code{hcfg}, \code{hcb}, \code{hwt}, \code{hsup}, \code{hprep},
\code{hrun}, \code{hwf}, \code{hbridge}.
\item \lead{Checked at every rung, not assumptions here}: \code{hcfg} (\code{Green.spike\_configPinned}),
\code{hwt} (\code{Witness.trExprS\_const\_of\_table}), \code{hsup} (\code{supportedB\_sound}),
\code{hwf} (\code{lbWfPeregrine\_of\_check} on a kernel-reduced Boolean, all twelve clauses of
\code{LBWfPeregrine}, \code{expandedFix} included).
\item \lead{Rung-only, universally quantified in the capstone}: \code{hev}, \code{hvwt},
\code{hty}, \code{hfo}; G5 additionally binds two fixture bundles and three defeq witnesses, G8
one fixture bundle and \code{hbody}.
\end{itemize}

\par\smallskip\noindent
\small
\begin{tabular}{lcccccccc}
\toprule
rung & G1 & G2 & G3 & G4 & G5 & G6 & G7 & G8 \\
\midrule
binders & 13 & 14 & 14 & 14 & 18\textsuperscript{*} & 14 & 14 & 16 \\
\bottomrule
\end{tabular}
\normalsize
\par\noindent
\textsuperscript{*}G5 proves \code{hev} rather than binding it.

\begin{itemize}
\item \lead{Caveat} two conclusion conjuncts are the theorem's own hypotheses --- \code{hprep}
and \code{hwf} --- so neither is established \emph{by} the theorem.
\item \lead{Caveat} \code{NoBodylessRefs} is no binder: each rung decides it by kernel
computation beside the theorem (\code{Green.g1\_noBodylessRefs} to \code{g8\_noBodylessRefs}).
\end{itemize}

\begin{assumption}[The reified source table is the live environment's slice]
\label{asm:htbl}
\uses{def:Witness-SourceTableAdequate, def:Witness-SourceTable}
\stAssumed{} \lead{In short} Every constant the reified table \code{tbl} pins is pinned as
\code{lenv} declares it, its tabled body is what \code{prepare\_erasure} computes up to
\hl{alpha-equivalence}, and every tabled inductive type is pinned.
\begin{itemize}
\item \lead{Given} three clauses. The body clause pins a tabled body only up to alpha-equivalence
(blind to binder names and binder info, nothing else) --- deliberately, because at equality the
clause is uninhabited wherever \code{inlineMatchers} draws \code{let} names from the name
generator.
\item \code{compilerLevels}: the \emph{compiler} declaration \code{getDeclInfo?} answers with ---
the \code{\_unsafe\_rec} companion where one exists --- carries the table's level column, which is
the scope \code{visitMutual} installs when it erases a member's body.
\item \lead{Class} D: no Lean term denotes the ambient \code{Lean.Environment}. \lead{Stands at}
the capstone and all eight rungs, discharged nowhere in Lean.
\item \lead{Mechanised outside}: \code{lake exe reify --check} compares the table against the
live environment, the level column included; green at all eight, with five G7/G8 bodies matching
only up to binder names and no mismatch.
\end{itemize}
\end{assumption}

\begin{assumption}[Every tabled declaration is safe, and the table's columns agree]
\label{asm:hsafe}
\uses{def:TableSafe}
\stAssumed{} \lead{In short} \code{TableSafe lenv tbl}, six clauses: every tabled constant,
inductive type and inductive-column constructor is \code{safe} in \code{lenv}, the table's
constant and constructor columns \hl{agree}, and every tabled body is \code{max}-free.
\begin{itemize}
\item \lead{Given} the three safety clauses feed Assumption~\ref{asm:ErasureSpec-decl_adequate}
(stated at safe declarations only); the two column clauses are \code{notUnsafeRec}, no tabled
constant is an \code{\_unsafe\_rec} companion, and \code{declCtor}, a tabled constructor is in the
constructor column too.
\item \code{noMaxLevels}: every tabled body lies in the \code{max}-free level fragment, which is
the fragment erasure transports along under level instantiation and what the bridge's level-column
premise spends.
\item \lead{Class} D; \lead{Stands at} the capstone and all eight rungs, discharged nowhere.
\code{lake exe reify --check} reads the live declarations for the three safety clauses (the safety
column itself is an upstream ask); the three column clauses are decidable on a concrete table.
\end{itemize}
\end{assumption}

\begin{assumption}[The blocks the run installs are tabled, lambda-headed and informative]
\label{asm:hblk}
\uses{def:TableBlocks}
\stAssumed{} \lead{In short} \code{TableBlocks lenv env tbl}, three clauses: every installed
block's members are tabled, every tabled body is a \hl{lambda}, and no member is erasable.
\begin{itemize}
\item \lead{Given} not a formality: \code{run\_mkDef\_box\_not\_lambda} exhibits an erasable
member with a non-lambda body; a blanket version would be \textbf{false} on every table (the
non-recursive instance constants and the rung subjects are non-lambda).
\item \lead{Class} D; \lead{Stands at} the capstone and all eight rungs, discharged nowhere ---
and \textbf{spent by no theorem on the capstone's chain}: \code{erasure\_bridge\_of\_run}
carries it unused.
\item \lead{Measured} over thirteen tables: 63 installed blocks, all singletons, no untabled
member, no non-lambda body.
\end{itemize}
\end{assumption}

\begin{assumption}[Every tabled compiler body is kernel-typeable at its declared type]
\label{asm:hcb}
\uses{def:CompilerBodies, def:lean4lean-grounding}
\stAssumed{} \lead{In short} For every constant the table gives a body for, it is declared in
\code{lenv} and the body \hl{translates and types} at the constant's declared type, read at the
constant's own level parameters.
\begin{itemize}
\item \lead{Stands at} the capstone and every rung; \textbf{checked at G1 only}
(\code{Green.g1\_compilerBodies}, \stChecked{}); binder at G2--G8.
\item \lead{Why it does not generalise}: G2--G4 carry \code{OfNat.ofNat} (needs \code{TrProj},
unproven at the pin); G5/G6 carry an eliminator spine and a recursive body; G7/G8's Arith table
carries ten projection-bearing bodies among thirty, twenty-eight of forty-four declarations
polymorphic.
\item \lead{Would be discharged by} running lean4lean's checker per body-type pair; what blocks
the ladder is \code{TrProj.weak'\_inv} and \code{TrProj.uniq}, \code{sorry} at the pin.
\end{itemize}
\end{assumption}

\begin{assumption}[The prepared term is the subject]
\label{asm:hprep}
\uses{def:Erasure-prepare_erasure, def:ConfigPinned}
\stAssumed{} \lead{In short} The preparation run on the subject succeeds from the empty state and
returns the prepared term with the state \hl{unchanged}; at a rung, the prepared term equals the
subject.
\begin{itemize}
\item \lead{Class} D: \code{IO.RealWorld} is opaque, so no Lean proof of a run's outcome exists.
\lead{Stands at} the capstone and all eight rungs, discharged nowhere in Lean; \code{lake exe
reify --prepared} reports identity, green at all eight.
\item \lead{Why it matters}: the capstone's syntactic conjuncts read the \emph{prepared} term while
\code{Supported}/\code{TrExprS} are checked at the subject; transport needs a stronger, generally
false, statement about \code{macroInline}.
\end{itemize}
\end{assumption}

\begin{assumption}[The \#erase run produced the committed program]
\label{asm:hrun}
\uses{def:Erasure-erase, def:ConfigPinned}
\stAssumed{} \lead{In short} \code{\#erase}, run on the subject at the pinned configuration,
returns the committed untyped environment and entry term, with a discarded list of inlining hints
(the \code{.ast.inlinings} sidecar).
\begin{itemize}
\item \lead{Class} D, the canonical permanent binder: \code{IO.RealWorld} is opaque. \lead{Stands
at} the capstone and all eight rungs, discharged nowhere; \code{lake exe green-check} re-runs
\code{\#erase} and byte-diffs the emitted \code{.ast} against the committed file.
\item \lead{Caveat} it is consumed through \code{erase\_run\_ok}, which splits the run into
preparation and \code{visitExpr} halves.
\end{itemize}
\end{assumption}

\begin{assumption}[The source evaluation of the subject at its spine]
\label{asm:hev}
\uses{def:SEval, def:SEvalFlags}
\stAssumed{} \lead{In short} The source semantics evaluates the subject applied to its spine to a
value, reading the compiler-body table. Class C: visible in the \hl{statement a reader checks}.
\begin{itemize}
\item \lead{Stands at} the capstone's observable clause and seven of the eight rungs, at the
empty spine (G8: one argument); \textbf{checked at G5 only} (\stChecked{}), via delta then iota at
\code{Nat.casesOn}.
\item \lead{Caveat} no witness at G7, G8: 45 recursive calls through \code{Nat.pow},
\code{Nat.mul}, \code{Nat.add}, \code{Nat.sub}, each owing a derivation, against G5's 45-line
precedent.
\item Volume, not vacuity: Lean is total, so every unselected branch of a well-typed closed term
has a value; \code{partial}/\code{unsafe} bodies are outside the fragment.
\end{itemize}
\end{assumption}

\begin{remark}
\lead{Caveat} \code{Green.g5\_seval} needs five further binders (two fixture bundles, three
defeq witnesses). The \code{Nat} bundle is satisfiable on its own
(Lemma~\ref{lem:Green-spikeNatFacts_natEnv}), but not together with \code{A}
(Section~\ref{sec:trust:vacuity}); the \code{PUnit} bundle has no witness.
\end{remark}

\begin{assumption}[The value's translation witness]
\label{asm:hvwt}
\uses{def:lean4lean-grounding}
\stAssumed{} \lead{In short} The value the source evaluation produced has a lean4lean translation
at the \hl{empty level scope and empty local context}. Class C, discharged nowhere.
\begin{itemize}
\item Its subject-side twin, \code{hwt}, \emph{is} checked at every rung (class A) but does not
reach the value: a rung's value is a constructor spine, not a constant.
\item \lead{Would be discharged by} a transport through \code{SEval.defeq}, moving the rung from
class A to class B and discharging Assumption~\ref{asm:hty} with it. No rung takes this route.
\end{itemize}
\end{assumption}

\begin{assumption}[The value is typed at the first-order inductive's spine]
\label{asm:hty}
\uses{def:lean4lean-grounding}
\stAssumed{} \lead{In short} The value's translation has a type headed by an inductive type
former --- at every rung, \hl{\code{Nat}}. Class C, discharged nowhere, for the reason
Assumption~\ref{asm:hvwt} is; this is what the first-order argument reads the answer's shape off.
\begin{itemize}
\item \lead{Would be discharged by} the same \code{SEval.defeq} route as
Assumption~\ref{asm:hvwt}.
\end{itemize}
\end{assumption}

\begin{assumption}[The answer's inductive is first order in the model]
\label{asm:hfo}
\uses{def:FirstOrderInd}
\stAssumed{} \lead{In short} \code{FirstOrderInd env I}, read at \code{Nat} by all eight rungs:
every constructor field of every block member is itself first order, and the block has \hl{no
parameter, no index, no \code{Prop} member}. Class C, discharged nowhere.
\begin{itemize}
\item \lead{Caveat} jointly \code{False} with \code{A} (Section~\ref{sec:trust:vacuity}), and so
is any discharge from \code{P}/\code{htbl}: \code{indInfo\_of\_tabled} alone contradicts \code{A}.
\item \lead{Blocked by} three facts: lean4lean's \code{TrEnv'.inductInfo\_inv} carries no
declaration list (the missing fusion is \code{TrEnv'.inductInfo\_wf'}); \code{TrEnv.inductInfo\_inv}
reads the kernel's \code{find?}, not \code{Lean.Environment.find?}; \code{BlockAdequate} is
per-name, so two members yield unrelated model blocks.
\item \lead{Caveat} the Boolean side \code{firstOrderIndB} is \code{true} at \code{Nat},
\code{Bool}, the fixture tree and \code{false} at \code{List}, \code{Prod}, \code{True}, but it
passes vacuously through \code{mdata} and \code{let}; no soundness theorem exists.
\end{itemize}
\end{assumption}

\begin{remark}
\lead{Caveat} The only proved \emph{model} inhabitant of \code{FirstOrderInd} in this tree is at
an empty inductive, so \code{firstorder\_erases\_deterministic} and \code{firstorder\_no\_box}
must never be presented as instantiated at \code{Nat}.
\end{remark}

\section{The measured axiom footprint}
\label{sec:trust:ledger}

\code{test/Ledger.lean} holds one \code{\#print axioms} invocation per tracked top-level result
--- 47 of them --- against the committed \code{test/ledger.expected} (394 lines);
\code{scripts/ledger.sh} re-measures and diffs, and CI fails on a mismatch.

\par\smallskip\noindent
\small
\begin{tabular}{p{0.46\linewidth}p{0.46\linewidth}}
\toprule
\textbf{theorems} & \textbf{measured footprint} \\
\midrule
\code{erases\_correct}, \code{erases\_correct\_lb}, \code{step\_delta}
& the four standard names plus the six-name instantiation cluster:
\code{Expr.mkAppData\_eq}, \code{Expr.mkData\_eq}, \code{Expr.replace\_eq},
\code{Level.hasParam\_eq} and the two \code{bv\_decide} certificates \\
\code{erases\_correct\_of\_steps}, \code{step\_iota}, \code{step\_proj},
\code{not\_erasable\_of\_informative}, \code{firstorder\_erases\_deterministic},
\code{firstorder\_no\_box}, \code{fOFields\_of\_asks}, \code{SEval.defeq}
& \code{propext}, \code{sorryAx}, \code{Classical.choice}, \code{Quot.sound} \\
\code{ErasesEnv.runtimeKey\_isCasesOn}, \code{erases\_elimSpine\_no\_value},
\code{ErasesEnv.ctorArity}, \code{CasesOnShape.agree}, \code{neverZeroB\_sound},
\code{lbWfPeregrine\_of\_check}, \code{constants\_of\_tabled},
\code{constOrigin\_of\_constants}, \code{TrExprS.alpha}, \code{Erases.alpha},
\code{SourceTableAdequate.erases\_prepared}, \code{visitExpr\_refines\_erasesLB},
\code{visitExpr\_refines\_erasesLBFix}, \code{bridgeEnv\_of\_regInv},
\code{bridgeEnv\_of\_regContent}, \code{Green.g5\_seval}, \code{Green.spikeNatFacts\_natEnv},
\code{Green.g8\_hargReach}
& \code{propext}, \code{Classical.choice}, \code{Quot.sound} \\
\code{ElimDecl.uniq}, \code{Lower.appReady}, \code{Lower.constToFix},
\code{LowerBlock.lambda\_of\_fixLambda}, \code{LBOptimize\_correct} & \code{propext}, \code{Quot.sound} \\
\code{lbEval\_sound} & \code{propext} \\
\code{Witness.Expr.AlphaEq.refl}, \code{Witness.Expr.AlphaEq.symm} & no axiom \\
\code{ErasureSpec.oracle\_sound\_of\_run},
\code{shipping\_erase\_correct\_firstorder}, \code{Green.green\_G1} through
\code{Green.green\_G8} & the identical 33-name set of
Assumption~\ref{asm:lean-reflection-axioms} \\
\bottomrule
\end{tabular}
\normalsize

\begin{enumerate}
\item The \code{sorryAx} in \code{erases\_correct} is \emph{inherited}, through uniqueness of
translations/definitional equality and through \code{VEnv.WF.patsStrong}; its three step arms
carry the same set.
\item \code{step\_proj} adds nothing to \code{step\_iota} (Assumption~\ref{asm:lean4lean-trust}).
\item The delta arm derives the instantiated reading of a tabled body's erasure itself, so it
carries the six-name instantiation cluster and hands it to \code{erases\_correct} and
\code{erases\_correct\_lb}. All six are already among the capstone's 33, so no rung's footprint
grows.
\item \code{not\_erasable\_of\_informative} cites \code{const\_arity\_inv} directly and
\code{fOFields\_of\_asks} reaches it through the Pi-spine peeling lemma; both rows carry
\code{sorryAx} through the injectivity cluster as well, so the citation moves no row.
\item \code{bridgeEnv\_of\_regInv} and \code{bridgeEnv\_of\_regContent} measure clean of both
clusters: three standard names, no \code{sorryAx}.
\item A second fixture, \code{test/erasesLB.expected}, measures the composite refinement's
introduction lemmas, statement \emph{and} footprint.
\end{enumerate}

\begin{remark}
\lead{Caveat} The graph's trust edges come from three sources: \emph{measured} (a ledger row),
\emph{read off the proof script} (the proof visibly spends a measured lemma), and the reflection
edge, drawn \emph{transitively through one theorem only} (capstone, rungs,
\code{ErasureSpec.oracle\_sound\_of\_run}). A node with neither edge is measured clean --- the
refinement theorems, \code{LBOptimize\_correct}, \code{lbEval\_sound}, \code{bridgeEnv\_of\_regInv}.
\end{remark}

\section{What the rungs and the benchmark programs actually cover}
\label{sec:trust:coverage}

Both tables are transcribed from \code{doc/coverage.md}. The five benchmark programs are frozen
copies under \code{VerifyBench/Src/}, held against the sibling benchmark repository; their roots
add the \code{\#erase} line setting \code{csimp := false}, which every correctness statement
needs and the frozen originals do not set.

\par\smallskip\noindent
\small
\begin{tabular}{llp{0.39\linewidth}p{0.26\linewidth}}
\toprule
\textbf{program} & \textbf{.ast bytes} & \textbf{in the fragment?} & \textbf{capstone} \\
\midrule
Arith & 14,241 & \textbf{yes} & applied capstone's subject, G7, G8 \\
Sieve & 28,555 & no --- recursor head at \code{Eq.rec} & not a rung \\
BinaryTrees & 30,056 & no --- recursor head & not a rung \\
Quicksort & 67,033 & no --- recursor head, well-founded division, sparse \code{casesOn} & not a rung \\
Fannkuch & 40,808 & no --- recursor head, eta-contracted minor & not a rung \\
\bottomrule
\end{tabular}
\normalsize

\hl{One of five} programs is inside the decidable fragment (Arith), the one the ladder
instantiates. Erroring-body counts: 0/43 (Arith), 2/81 (Sieve), 2/89 (BinaryTrees), 8/126
(Quicksort), 5/95 (Fannkuch). \code{NoBodylessRefs} holds on all five: the registering exits give
a reachable recursor or quotient primitive a body, so no emitted program carries a body-less
declaration.

\begin{itemize}
\item \lead{Not exercised by any of the five}: a genuinely mutual fixpoint block (all 63 installed
blocks across the thirteen measured tables are singletons), \code{Acc}, \code{WellFounded},
\code{Quot}, and any elimination with a \code{Prop} discriminee --- the fragment excludes such
eliminations outright.
\end{itemize}

G1--G7 are closed nullary definitions ([L]'s Theorem 15); G8 is the Arith function applied to its
argument, the one rung with a non-empty spine. Each rung's conclusion ends in a literal peano
numeral, so no box or stuck term can satisfy it.

\par\smallskip\noindent
\small
\begin{tabular}{llllrllll}
\toprule
\textbf{rung} & \textbf{subject} & \textbf{answer} & \textbf{.ast bytes} & &
\textbf{hcb} & \textbf{hev} & \textbf{hbridge} & \textbf{hbody} \\
\midrule
G1 & \code{spikeZero} & 0 & 354 & & checked & binder & binder & --- \\
G2 & \code{spikeLit} & 4 & 1,556 & & binder & binder & binder & --- \\
G3 & \code{spikeLet} & 3 & 1,469 & & binder & binder & binder & --- \\
G4 & \code{spikeProj} & 1 & 2,284 & & binder & binder & binder & --- \\
G5 & \code{spikeCase} & 1 & 966 & & binder & \textbf{checked} & binder & --- \\
G6 & \code{spikeFix} & 2 & 885 & & binder & binder & binder & --- \\
G7 & \code{arithClosed} & 8 & 14,625 & & binder & binder & binder & --- \\
G8 & \code{benchArith} at \code{Nat.zero} & 8 & 14,291 & & binder & binder & binder &
\textbf{binder} \\
\bottomrule
\end{tabular}
\normalsize

\begin{itemize}
\item G6's subject is applied to \code{Nat.succ (Nat.succ Nat.zero)} in constructor form, to
exercise delta on a recursive constant applied to a \emph{constructor} value.
\item \lead{Settled by computation at every rung}: \code{hcfg}, \code{hsup}, \code{hwf},
\code{NoBodylessRefs}, the target-side evaluation, and \code{hwt} by a checked term. The
elaboration environment and its model stay universally quantified.
\item \lead{Caveat} \hl{no instantiation satisfies them}: \code{A} is jointly \code{False} with
\code{hfo} and with \code{P}/\code{htbl} at every rung (Section~\ref{sec:trust:vacuity}), so a
literal answer shows only that the checked terms meet the stated interface.
\item \lead{Non-vacuity guards} \code{LBWfPeregrine}'s binder-name clause is a printability test,
with offending-name count 0 at all eight rungs; \code{test/Vacuity.lean} holds the level-scope
regression and the refutations of Section~\ref{sec:trust:vacuity}.
\item \lead{Dead-declaration budget}: \code{lake exe hygiene --dead} reports 343 declarations
outside the import closure --- tooling (107), frozen benchmark sources (43), the optimiser (71),
the uniform-erasure module (19), \code{test/Vacuity.lean} (9), the twelve \code{test/fixes/}
regressions (94) --- exception list empty.
\end{itemize}

\section{Shipping findings}
\label{sec:trust:findings}

Twelve findings against the shipping eraser carry a fix, a guard or a derived fact in the code,
each with a regression test under \code{test/fixes/} run by \code{scripts/fixes.sh}; one more,
F-PRODUCT, is an owner decision. \code{doc/rework/03-DEV-FIX.md}'s applied-edits table is the
index. The column says what the code guarantees.

\par\smallskip\noindent
\small
\begin{tabular}{lp{0.53\linewidth}p{0.26\linewidth}}
\toprule
\textbf{id} & \textbf{what the code guarantees or guards} & \textbf{regression test} \\
\midrule
F-PROP & \code{register\_inductive} sets \code{propositional} to the arity decision on the
inductive's own type, and the field has no default, so it cannot be omitted &
\code{test/fixes/F-PROP.lean} \\
F-ETA & the registration loop registers the eta-expanded fixpoint, so every emitted fixpoint node
sits under an argument spine and \code{LBWfPeregrine.expandedFix} holds &
\code{test/fixes/F-ETA.lean} \\
F-ETA2 & the eta loops bind each already-supplied non-value argument in a \code{let} outside the
binders they open, so no supplied argument is re-run per application &
\code{test/fixes/F-ETA2.lean} \\
F-SPARSE & \code{visitCases} reads the inductive from \code{casesInfo.indName}, expands a
catch-all into one alternative per uncovered constructor, and refuses the shapes it cannot compile
with an error instead of an assertion & \code{test/fixes/F-SPARSE.lean} \\
F-EQREC & a recursor with no compiler value is registered with \code{recursorRealizer}'s
synthesized body rather than as a body-less axiom & \code{test/fixes/F-EQREC.lean} \\
F-QUOT & a quotient primitive is registered with \code{quotRealizer}'s body at the arity the
kernel fixes & \code{test/fixes/F-QUOT.lean} \\
F-ACC & \code{visitCases} refuses an elimination of a propositional inductive one of whose
constructor fields is not a proof & \code{test/fixes/F-ACC.lean} \\
F-DEPTH & \code{isArityCheck} fuels its walk from a fixed budget independent of the unreduced
subject's syntactic depth & \code{test/fixes/F-DEPTH.lean} \\
F-UNSAFEREC & \code{visitMutual} refuses a block whose mapped kernames repeat, and
\code{run\_rec\_exit\_nodup} carries that distinctness out of a successful run &
\code{test/fixes/F-UNSAFEREC.lean} \\
F-KERNAME & \code{checkKernameFresh} and \code{checkIndKernameFresh} refuse a key already minted
for another name at every registration point; the encoding stands and
\code{toKername\_not\_injective} (Lemma~\ref{lem:toKername_not_injective}) remains true &
\code{test/fixes/F-KERNAME.lean} \\
F-DEPLCTX & \code{visitMutual}'s two dependency re-entries reset \code{ErasureContext.lctx} to
empty alongside the level column, so a dependency's closed body is erased in the empty context,
as MetaRocq's \code{erase\_constant\_body} does; emitted bytes are unchanged on 21 programs &
\code{test/fixes/F-DEPLCTX.lean} \\
F-ARITYLET & \code{arityResultSort} walks \code{.letE} and \code{.mdata}, as \code{destArity}
walks \code{tLetIn}, and does not reduce, so an arity behind a constant stays non-propositional &
\code{test/fixes/F-ARITYLET.lean} \\
F-PRODUCT & \code{auto\_inline\_typeclass\_dispatch}, an unverified product feature, default off
and pinned \code{false} by \code{ConfigPinned}; an owner decision, class E & --- \\
\bottomrule
\end{tabular}
\normalsize

\begin{itemize}
\item \hl{No fix widens the verified fragment}: the scope restrictions stand, the first-order
observable stands, and Arith is the one benchmark program inside the fragment. Each fix turns a
silent miscompile or a stuck program into a correct one or an explicit refusal.
\item \lead{Guard} \code{isArityCheck}'s loop throws on fuel exhaustion (F-FUEL), routing the
verdict to the \code{Meta} fallback instead of answering \code{false}.
\item \lead{Caveat} sixteen panic sites --- eleven unreachable-assertions, five explicit ---
print to \code{stderr} and return the \code{Inhabited} default, which in the erasure monad is
$\Box$, so \textbf{a panic does not stop the run}. No predicate asserts a run did not panic: the
shape induction is panic-\emph{tolerant}, each site excluded by a proved invariant.
\end{itemize}

\section{What is open}
\label{sec:trust:open}

Ordered as \code{doc/rework/07-STATUS.md} section 4 orders them: the first item gates the rest.

\begin{enumerate}
\item \lead{The whole-ladder vacuity} \stOpen{} Read \code{ConstOrigin}, \code{CtorOf},
\code{IndInfo}, \code{CasesOnShape} and \code{IndBlockBelow} off \code{env}'s own declaration
list (\code{env.WF' ds}), as MetaRocq's \code{lookup\_env} does. Exclusion and uniqueness then
follow from lean4lean's \code{consts\_origin} with no fork change, \code{UpstreamAsks} shrinks
to at most \code{mkAppsInv} and \code{indSpineInj}, and the rungs stop being vacuous. The
\code{mono} lemmas do not survive unchanged: a list is not monotone in the environment order.
\item \lead{\code{hbridge}} \stOpen{} The accumulator bundle (\code{AccState}, \code{AccGrows},
eighteen motives, \code{AccAsks}; Section~\ref{sec:coldstart:acc}) has seventeen of eighteen steps proved; \code{StepAcc6},
\code{visitMutual}'s step, is open and needs the block-exit decomposition, an eighteen-motive
bundle for \code{ConstsDeclared} at the emitted term, the \code{Erases} induction behind it, the
constant-key \code{SpecEntryOk} layer and the block exit's prefix from a run.
\code{AccGrows} has no producer and \code{motives\_of\_steps\_acc} no consumer. Then
\code{erasure\_bridge\_env} retires \code{hbridge} (\code{doc/rework/12-REPAIRS-W9.md}).
Five proved steps (3, 10, 13, 14, 17) follow from \code{AccAsks.upstream} alone and need
re-proof after item 1.
\item \lead{\code{hfo}} a round-5 lean4lean commission amended with \code{TrEnv'.inductInfo\_wf'},
and a specification decision on the two \code{find?} primitives (Assumption~\ref{asm:hfo}).
\item \lead{Unspent binders} \code{consts\_classified} and \code{constOrigin\_of\_constants}
carry an unused \code{A}.
\item \lead{Upstream asks 9 and 10} and the two \code{TrProj} \code{sorry}s, drafted for round 5;
the \code{TrProj} pair blocks Assumption~\ref{asm:hcb} at G2--G8: ten of G7's 30 tabled bodies
carry a class projection.
\item \lead{Standing} \code{hev} constructed at G5 only; \code{TableSafe.noMaxLevels}; the alpha
gap (five G7/G8 bodies match up to binder names); $\Lower$ is not a function at a block member;
\code{SchemeNames} has no producer; four of five programs are outside the fragment; F-PRODUCT.
\item \lead{Consumers}: peregrine's Lean tests and the frontend benchmark pin the default branch,
and \textbf{neither builds the verified eraser}.
\end{enumerate}

\begin{remark}
\lead{Caveat} \code{doc/rework/12-REPAIRS-W9.md} and \code{11-REPAIRS-W8.md} are plans: read them
for the shape of remaining work, not for status. \code{bridgeEnv\_of\_regContent}'s docstring
says no theorem produces \code{RegAcc} or \code{RegKeyed} at a run; \code{regKeyed\_of\_run} and
the two registration producers contradict it in part.
\end{remark}

\noindent\lead{The census: what stands after \code{hbridge}} Closing it moves one row; item 1
removes \code{constOriginExcludes}.

\par\smallskip\noindent
\small
\begin{tabular}{lp{0.84\linewidth}}
\toprule
\textbf{class} & \textbf{what stands} \\
\midrule
D & \code{ErasureSpec}, seven fields no Lean term can state agreement with, because
\code{Lean.Environment} and \code{IO.RealWorld} are opaque; \code{htbl}, \code{hsafe},
\code{hblk}, \code{hprep}, \code{hrun}, mechanised outside Lean \\
C & \code{EraserAsks}, four fields; \code{hcb} at G2--G8; the per-rung binders \code{hev},
\code{hvwt}, \code{hty}, \code{hfo}; \code{hbody} at G8 \\
upstream & \code{UpstreamAsks}, four fields, one refuted; the seventeen lean4lean \code{sorry}
roots \\
scope & the fragment predicate, the \code{max}-free level fragment, the first-order
forward-simulation shape of the observable, and everything downstream of the emitted program
value \\
\bottomrule
\end{tabular}
\normalsize

\section*{Deviations from the paper, and caveats}

\begin{itemize}
\item The capstone is conditional: fourteen binders, three bundles. At the eight rungs, four
(five at G1) become checked terms; the rest, including \code{hbridge}, stand assumed.
\item No environment satisfies a rung's hypotheses: \code{UpstreamAsks} is refuted wherever an
inductive block is declared, so every rung is vacuous until the declaration-list restatement.
\item Two \code{EraserAsks} fields are not theorems and are not expected to become ones without
upstream work: \code{kernel\_ind\_head\_true} leaves a residue no constant budget bounds, and
\code{passes\_sound} quantifies a spine no pass respects syntactically.
\item The default \code{\#erase} configuration is \textbf{not} the pinned one: it differs from
\code{ConfigPinned} on three of five fields, so a default invocation is outside every theorem here.
\item The step from the emitted program value to the \code{.ast} bytes peregrine reads is outside
every theorem: the serialiser is a \code{partial def} with no equational lemma.
\item Departures from [JACM]/[L]: applied-form constructors, the eta-expanded fixpoint wrapper the
registration emits, Lean-specific $\Lower$ steps, and the first-order restriction.
\end{itemize}
